\section{Conclusion} \label{conclusion}
\vspace{-1mm}
We introduced \texttt{BiMol-Diff}, a unified diffusion framework addressing the bidirectional tasks of molecular graph generation and textual captioning. By moving beyond standard data-agnostic corruption to a \textit{token-aware noise schedule}, our approach explicitly preserves chemically salient substructures during the diffusion process. Empirical validation on M3-20M and ChEBI-20 confirms that this strategy yields significant gains over state-of-the-art autoregressive and diffusion baselines, particularly in structural fidelity (S2G) and semantic alignment (G2S). These findings demonstrate that token-aware denoising is essential for high-fidelity, controllable molecular language modeling, offering a robust alternative to autoregressive approaches in scientific domains.

\section{Limitations}
First, \texttt{BiMol-Diff} relies on linearized sequence representations and standard attention, lacking the explicit structural inductive bias inherent to graphs, which may limit data efficiency for complex topologies. Second, as with most diffusion models, our approach incurs higher computational costs during inference compared to autoregressive baselines due to the iterative denoising process, highlighting the need for future work on accelerated sampling.